% !TeX root = ../深度学习与大语言模型技术原理.tex
% 由主文档通过 \input 载入；请编译主文档。
\setcounter{chapter}{13} % 保留原书章号。
\chapter{稀疏检索、稠密检索与重排序}
\label{ch:retrieval}

\section{倒排索引与 BM25}

检索任务给定查询 $q$ 和文档库 $\mathcal C$，输出与查询相关的候选文档。稀疏检索（sparse retrieval）主要通过词项匹配；倒排索引（inverted index）为每个词项存储包含它的文档及频率等信息，避免对所有文档逐一扫描。TF--IDF 用词频和逆文档频率赋权，使常见功能词的影响较小，但不同实现的词频、平滑和归一化约定有所不同。

BM25 的一种常用形式为
\begin{equation}
\operatorname{score}(q,d)=\sum_{t\in q}\operatorname{IDF}(t)
\frac{f(t,d)(k_1+1)}{f(t,d)+k_1(1-b+b|d|/\overline\ell)},
\label{eq:bm25}
\end{equation}
其中求和按查询中的不同词项进行，$f(t,d)$ 是词频，$|d|$ 为文档长度，$\overline\ell$ 为库中平均长度，$k_1>0$ 控制词频饱和，$b\in[0,1]$ 控制长度归一化。可取 $\operatorname{IDF}(t)=\log(1+(N-n_t+0.5)/(n_t+0.5))$，$N$ 为文档数，$n_t$ 为包含该词的文档数\cite{bm25}。

词频分式随 $f$ 增大趋向 $k_1+1$，避免一个词重复很多次就无限提高相关性。长文档的分母增大，部分补偿长文天然更易匹配词项的优势。标识符、错误码、精确数字和人名往往需要这种词项匹配；同义表达和跨语言匹配则可能受限。分词与大小写处理同样会影响结果。

\section{双编码器与对比学习}

稠密检索（dense retrieval）用查询编码器 $f_\theta(q)$ 和文档编码器 $g_\theta(d)$ 映到向量空间，称为双编码器（bi-encoder）。得分可取点积 $s(q,d)=f_\theta(q)^\top g_\theta(d)$；若都做单位范数归一化，点积就是余弦相似度。未归一化时，向量范数也参与排序，不能将二者混称\cite{dpr}。

给定正文档 $d_i^+$ 和负例集合 $\mathcal N_i$，对比损失可写为
\begin{equation}
\ell_i=-\log\frac{\exp(s(q_i,d_i^+)/\tau)}{
\exp(s(q_i,d_i^+)/\tau)+\sum_{d\in\mathcal N_i}\exp(s(q_i,d)/\tau)},\qquad \tau>0.
\label{eq:contrastive}
\end{equation}
批内负例（in-batch negatives）复用同批其他查询的正例文档，降低额外编码成本；困难负例（hard negatives）选择词面或语义接近但不满足查询的文档，提高区分能力。若负例其实相关，就会错误推开应靠近的表示，因此必须检查假负例率。

编码器更新后，旧索引中的向量可能与新查询向量不再处于同一空间。发布检索模型应同步管理编码器版本、分词器、向量归一化、文档版本和索引版本。只更新查询端权重而不重建文档向量，除非方法特意支持，否则会产生难以诊断的退化。

\section{近似最近邻与索引取舍}

精确向量搜索对 $N$ 个 $d$ 维向量逐一计算得分，单查询约需 $O(Nd)$ 运算。近似最近邻（approximate nearest neighbor, ANN）通过索引减少访问量。倒排文件（inverted file, IVF）先将向量聚类，再只搜索部分簇；层次可导航小世界图（HNSW）通过图上的局部跳转接近近邻；乘积量化（product quantization, PQ）把向量分成子空间，各子空间以码本索引近似存储。

这些机制分别减少搜索范围、利用邻接结构或压缩向量。簇探测数、图搜索宽度和码本精度越低，通常延迟与内存更小，但 ANN 召回可能下降。应相对精确向量搜索测量索引损失，再把表示模型本身的语义召回单独测量，否则无法分清是向量不好还是索引遗漏。

FAISS 提供向量检索与索引算法，Elasticsearch、OpenSearch 组织文本索引及相关检索服务，Milvus、Qdrant 等向量数据库还处理持久化、元数据过滤和服务管理。选择工具时应从数据更新、过滤、规模与延迟要求出发；数据库品牌本身不决定检索质量。

\section{混合召回、RRF 与分数融合}

混合检索（hybrid retrieval）合并词项与向量召回，以利用互补错误。若不同检索器得分量纲不同，直接相加可能让一个检索器支配结果。加权分数融合需要校准或归一化；归一化也不能保证不同查询下分数含义一致。

倒数排名融合（reciprocal rank fusion, RRF）只使用排名：
\[
\operatorname{RRF}(d)=\sum_{j:d\in\mathcal C_j}\frac{w_j}{c+\operatorname{rank}_j(d)},\qquad c>0.
\]
未被检索器 $j$ 召回的文档不贡献该项。$c$ 控制对头部排名的敏感性，$w_j$ 控制各路权重。它减少对分数标度的依赖，但丢弃了得分间距信息，也不能修复所有召回器共同漏掉的文档。

\section{交叉编码器与迟交互}

重排序器（reranker）在较小候选集上进行更昂贵的相关性建模。交叉编码器（cross-encoder）将查询和文档拼接输入同一网络，使每层都可建模二者交互，再输出相关分数；因为文档表示依赖查询，难以像双编码器那样离线一次编码后直接复用。候选数量与文档长度直接影响重排成本。

迟交互（late interaction）保留词元级表示，在最后阶段计算细粒度匹配。例如 ColBERT 风格得分为
\[
s(q,d)=\sum_{i=1}^{|q|}\max_{1\le j\le|d|}q_i^\top d_j,
\]
其中向量通常归一化，并对有效词元使用掩码。它让每个查询词元寻找最强文档匹配，保留更多局部细节，但每文档需要存储多个向量\cite{colbert}。召回、重排和最终生成应分层测量；重排不能从不存在于候选集的文档中创造证据。

